\paragraph{Full-length comparison at a fixed gap.}
The short-domain condition is a substantive distinction. The following
elementary comparison already rules out a list bound depending only on
rate and gap when the domain occupies almost the whole prime field.
We make no novelty claim for this full-length observation.

\begin{proposition}\label{prop:full-length-fixed-gap}
Let $p\equiv1\pmod8$ be prime, $n=p-1$, and $k=n/4$.
On $\F_p^\times$, there is a received word with $n/2$ distinct
degree-$(k-1)$ polynomials, each agreeing on exactly $3n/8$ coordinates.
Thus the rate is $1/4$ and the gap is $1/8$. For $p>256$, the radius
$5/8$ is strictly below the characteristic-based Elias radius.
\end{proposition}
\begin{proof}
Let $\chi$ be the quadratic character, with $\chi(0)=0$, and put
$e=(p+1)/2=2k+1$. For $a\ne0$ define
\[
 G_a(X)=\sum_{j=0}^k\binom e{2j+1}a^{e-2j-1}X^j,
 \qquad P_a(X)=G_a(X)-X^k,
 \qquad w(x)=\frac{1+\chi(x)}2-x^k.
\]
The degree-$k$ terms cancel, and
$[X^{k-1}]P_a=-a^2/8$. All powers of $a$ are even, so the parameters
$a$ modulo sign give exactly $n/2$ distinct polynomials of degree $k-1$.
The binomial identity
\[
 G_a(s^2)=\frac{(a+s)^e-(a-s)^e}{2s}
\]
counts their agreements. If $x=s^2$ is nonsquare, take $s$ in
$\F_{p^2}$. Then $s^p=-s$, and $u=(a+s)^e$ satisfies
$u^2=a^2-x$ and $(a-s)^e=u^p$. Hence $G_a(x)=0$ exactly when
$\chi(a^2-x)=1$. There are $(p-1)/4$ such nonsquare $x$.

For square $x=s^2$ with $s\in\F_p^\times$, Euler's criterion gives
\[
 G_a(s^2)=
 \frac{(a+s)\chi(a+s)-(a-s)\chi(a-s)}{2s}.
\]
Away from $s=\pm a$, this equals one exactly when both characters
are one. The weighted count
$\frac14\sum_s(1+\chi(a+s))(1+\chi(a-s))=(p-1)/4$
uses $\sum_s\chi(a^2-s^2)=-1$. Removing $s=0$ and correcting the
half-weights at $s=\pm a$ cancel: at the endpoints the actual value is
$\chi(2a)=\chi(a)$. Dividing by two for $s$ and $-s$ gives
$(p-1)/8$ square agreements. The nonsquare count above follows from
the same quadratic-character identity, applied to
$(1-\chi(x))(1+\chi(a^2-x))/4$ with its endpoints removed.
Thus the total is $3n/8$. Finally,
$H_p(5/8)\le5/8+1/\log_2p<3/4$ for $p>256$.
\end{proof}

Here $n/p\to1$. In contrast, Theorem~\ref{im:all-list} and
Corollary~\ref{gm:gap-field} give exponentially large lists on short
integer intervals, with the stated quantitative dependence on the gap.
The full-length example does not settle the fixed-gap, short-domain
proximity-gap question. Exact checks of the displayed agreement counts
are in \path{research/dickson_fixed_gap/}.

\begin{lemma}[Splitting of pairwise differences]\label{lem:dickson-difference-splitting}
For the polynomials above, if $a^2\ne b^2$, every root of $G_a-G_b$
belongs to $\F_p$.
\end{lemma}
\begin{proof}
Let $t^2$ be a root in an algebraic closure, with
$y=G_a(t^2)=G_b(t^2)$. The case $t=0$ is immediate. For $c\in\{a,b\}$,
put $U_c=(c+t)^e$ and $V_c=(c-t)^e$. Then
\[
 U_c-V_c=2ty,\qquad U_c^2-V_c^2=2c(t+t^p),\qquad
 U_c^2+V_c^2=2c^2+2t^{p+1}.
\]
If $y=0$, then $t^p=-t$, so $t^2\in\F_p$. Otherwise put
$H=(t+t^p)/(ty)$. We have $U_c+V_c=cH$, whence
\[
 c^2H^2+4t^2y^2=4c^2+4t^{p+1}.
\]
Subtracting the equations for $a$ and $b$ gives $H^2=4$, and then
$t^2y^2=t^{p+1}$. Thus $(t+t^p)^2=4t^{p+1}$, or
$(t^p-t)^2=0$. Again $t^2\in\F_p$.
\end{proof}

\begin{corollary}[Quadratic exceptions in growing prime characteristic]
\label{cor:quadratic-extension-fixed-gap}
For every prime $p\equiv1\pmod8$, $p\ge17$, put $n=2(p-1)$.
Over $\F_{p^2}$ there is a Reed--Solomon code of length $n$ and dimension
$n/8$, and a received line with at least $\lceil n^2/20\rceil$
full agreement-set MCA exceptional labels at agreement $3n/16$.
Thus the capacity gap is exactly $1/16$, and $p=n/2+1>n/8-1$.
\end{corollary}
\begin{proof}
Set $N=p-1$. Proposition~\ref{prop:full-length-fixed-gap} supplies
$L=N/2$ candidates at agreement $A=3N/8$, with dimension $k=N/4$.
Choose any nonsquare anchor. Exactly $\ell=N/4$ candidates agree there:
indeed $\sum_{a\in\F_p}\chi(a^2-x)=-1$ for nonsquare $x$, none of
the summands is zero, and the $a=0$ summand is $-1$. There are therefore
$N/2$ admissible nonzero parameters, or $N/4$ modulo sign.
As in Proposition~\ref{prop:exact-list-amplification}, divide each anchored
candidate minus the received anchor value by the anchor's linear factor.
The resulting quotients have degree at most $k-2$. By
Lemma~\ref{lem:dickson-difference-splitting}, their evaluations are pairwise
distinct at every point outside $\F_p$.

Append any $N+1=p$ distinct points of $\F_{p^2}\setminus\F_p$, with
direction one, and retain direction zero on the $N-1$ old coordinates.
Independent uniform translations of the padding value sets yield an
expected number of counted labels equal to
\[
 p^2\bigl[1-(1-\ell/p^2)^p\bigr]
 \ge p^2\bigl[1-\exp(-N/(4p))\bigr]
 \ge Np\bigl[1-\exp(-1/4)\bigr]
 >N^2/5=n^2/20.
\]
The second inequality uses concavity of $1-\exp(-u/4)$ on $[0,1]$;
the last uses $1-e^{-1/4}\ge1/4-1/32=7/32>1/5$.
Every counted label has a selected quotient with at least $A-1\ge k$
old agreements and one new agreement. A degree-$<k$ direction polynomial
matching this full support would vanish on the old agreements, hence
identically, but would equal one on the new agreement. This is impossible.
The dimension and agreement threshold remain $k$ and $A$.
\end{proof}

This elementary consequence rules out a universal linear capacity MCA
bound in the large-characteristic field class of~\cite{dkt2026}.
Its ambient field is a quadratic extension, not a prime field, and its
agreement lies below the first-order regime. It therefore neither proves
tightness of the first-order quadratic upper bound nor resolves the
prime-field superlinear-exception target. It concerns full agreement-set
MCA; absence of ordinary correlated agreement is not asserted.
Exact extension-field replays at $p=17,41$ are recorded in
\path{research/prime_field_tightness/quadratic_extension_verification.json}.
